\section{Conclusions}
\label{paper:conclusiones}
L’extrême et moyenne raison holographique se trouve formulée au moyen d’une structure génératrice, d’une évaluation arithmétique réversible et d’une loi de transport avec récupération d’information. Le registre dodécaphasique \(K\), produit par APP--TRIT--TPK, conserve son ordre dans le système de coordonnées de phase et relie la clôture arithmétique à la représentation incidentielle. Sa lecture
\[
\kappa=
\frac{234543140729659824621058914794146601}{10^{36}-1}
\]
récupère exactement les douze composantes : la multiplication par le dénominateur restitue l'entier et la base mille détermine ses blocs. Ce résultat fixe le sens mathématique de la constante holographique de couplage arithmético--géométrique. Sa valeur est une représentation exacte du registre ordonné et participe, avec les coordonnées régionales de \(\pi,\varphi,e\), à la détermination de \(\alpha_A\).

La même structure permet d’identifier des opérations. Les douze décalages cycliques de \(K\) forment une base dont la matrice de Gram a pour extrémités spectrales \(589609\) et \(6263^2\). Les réponses à cette base déterminent un opérateur linéaire de facteur d’amplification de l’erreur \(1/\sqrt{589609}\) ; une réponse vectorielle suffit dans la classe qui commute avec le décalage. La représentation incidencielle transporte le registre vers son image spécifiée avec un inverse explicite. Deux observations de mémoire retrouvent sa partie centrée avec la borne optimale \(9/4\), et la charge uniforme complète le registre. Évaluation, incidence et identification se composent ainsi sur un même objet, en conservant à chaque étape les données requises par l’inverse.

La coïncidence des représentations entière et analytique de \(\alpha\) est établie à toute profondeur. Le système de coordonnées analytique prolonge la troncature d’ordre neuf au moyen d’une récurrence de raison \(1/729\), dont le coefficient initial est obtenu à partir de la précoordonnée déjà engendrée. La convergence uniforme, la borne positive de la dérivée et les intervalles rationnels emboîtés identifient la même racine simple et les mêmes préfixes canoniques dans les deux représentations. Le théorème détermine leur compatibilité complète au sein de cette construction corrélée.

Le résultat général qui soutient la conservation est l'identité bilatérale
\[
 (a+1)(aB-C)^*(aB-C)
 +a\bigl((a+1)B+C\bigr)^*\bigl((a+1)B+C\bigr)
 =\bigl((a+1)^3+a^3\bigr)I,
\]
pour \(a>0\) et deux isométries \(B,C\) avec domaine et codomaine communs. L'annulation des termes croisés démontre l'égalité en dimension arbitraire. Dans la spécialisation \(a=8\), le facteur ternaire \(73\) conserve son interprétation comme norme algébrique et indice réticulaire, et satisfait \(10\cdot73=729+1\). Le raffinement d'étiquettes entières admissibles possède un inverse par division euclidienne, avec profondeur et convention de retenue spécifiées ; son relèvement aux bases d'états est unitaire et se compose avec le bilan.

La loi permet également de déterminer ce qui change lors de l’observation. Pour un transport relatif unitaire \(R\), le lecteur diagonal transforme l’observable du système de coordonnées commun selon
\[
 \mathcal T(G)=\frac{72}{73}G+\frac1{73}R^*GR.
\]
Ses points fixes sont exactement les observables qui commutent avec \(R\). L’observable transporté sur l’état complet conserve, en revanche, la lecture initiale par sa formule de restitution. Les identités extérieures prolongent ce résultat à chaque degré admissible, en conservant les secteurs purement observés, les secteurs enregistrés et les secteurs mixtes. Les contributions d’aire et de volume sont incluses dans une loi de formes, avec les métriques et les dimensions correspondantes.

La récupération est maintenue à une profondeur arbitraire au moyen d’une politique explicite d’archivage. Dans le cas nonadique, l’opérateur terminal satisfait \(\|R_N\|^2\leq(729/1241)^N\). L’archive infinie est donc une isométrie et son adjoint reconstruit l’état avec une norme opératorielle un. Pour chaque profondeur positive, l’archive finie admet un inverse exact corrigé par son gramien, avec pour borne
\[
 \|L_N\|\leq\bigl(1-(729/1241)^N\bigr)^{-1/2}.
\]
Ces formules quantifient la stabilité et séparent l'inversion exacte de l'approximation obtenue par troncature d'une série de récupération. Dans les dynamiques générales, la mémoire et le terminal survivant conservent conjointement l'état ; la politique indiquée démontre le régime dans lequel toute sa norme est transférée à l'archive.

Les réalisations physiques développées appliquent ces opérations avec leurs données explicites d’action, d’orientation, d’échelle et de représentation. En particulier, pour le protocole symplectique déclaré avec deux entrées gaussiennes pures, centrées et identiques et un nombre fini de modes, la composante antilinéaire de la mémoire fixe le spectre symplectique de la réduction, en maintenant toutes les occupations de Fock. La boucle elliptique--parabolique produit la pureté \(9/11\), le spectre \((9/10)10^{-j}\) et l’entropie adimensionnelle \((10/9)\log10-\log9\). La conservation globale et le mélange de l’observation réduite sont décrits dans le même protocole. Les réalisations électromagnétiques, gravitationnelles, thermiques et variationnelles maintiennent les opérateurs et les normalisations de leurs preuves respectives.

La construction commune du continu conserve les cinq opérations corrélatives et la compatibilité de leurs raffinements. La clôture des collections d’histoires de la section~\ref{sec:cierre-cardinal-nativo} établit, dans l’univers généalogique qui y est défini, l’égalité cardinale interne de sa droite avec \(\aleph_1\) de cet univers. L’affirmation conserve le domaine et les opérations d’appartenance qui interviennent dans ses deux injections et dans son identification finale.

Le recensement de l'émetteur résiduel explicite la base finie d'une partie de cette chaîne : il réunit \(44\) signatures, \(18\) additives et \(26\) multiplicatives, avec leurs multiplicités et la règle de reconstruction des fibres. Les \(648\) appartenances de conditions initiales, \(324\) dans chaque feuillet, sont conservées dans les données du supplément. Les sept théorèmes vérifiés dans Lean certifient les identités arithmétiques et leurs conséquences finies déclarées ; les démonstrations d'opérateurs, de limites et d'incidences sont développées dans le texte. Le recensement, le calcul et la formalisation fixent ainsi leurs objets de vérification respectifs.

% BEGIN RECIPROCIDAD 20260918

L’extension des sections~\ref{rec:positiva}--\ref{rec:recuperacion} ajoute une réalisation géométrique récupérable de cette continuité. La direction de \(K\) détermine un flot positif unimodulaire, un couple de réseaux duaux et une géodésique à rayons compensés. L’archive reconstruit la direction et contrôle l’erreur de ses lectures thêta et Epstein. La structure polaire demeure constante, tandis que la variation spectrale est entière. Par une autre voie de la même généalogie, le contraste angulaire est réalisé dans le système de coordonnées de Klein et retrouve la section de Planck avec son orientation ; sa substitution dans l’expression électronique conserve tous ses facteurs. La séparation de l’opérateur de masses distingue l’échelle commune de la congruence de déterminant un. Ces compositions élargissent la fonction opératoire de la constante de couplage sans identifier les différents domaines de conservation.
% END RECIPROCIDAD 20260918

\newpage % REV03: balance complete concluding blocks.
% BEGIN CENTRALIDAD_K_20260921
La sélection conjointe de la constante est développée avant sa
reconstruction algébrique : les ordonnancements du répertoire terminal produisent
\(19\,446\) registres ; la compatibilité exceptionnelle en conserve \(79\), et la
compatibilité temporelle détermine un unique registre \(K\). On conserve les
conditions des deux filtres et leur dépendance à l’égard de l’histoire enrichie.
La section~\ref{act21:centralidad-constante} articule cette détermination avec
la clôture arithmétique, les représentations incidencielles et la récupération
opératoire.
% END CENTRALIDAD_K_20260921

% BEGIN SINTESIS_ADITIVA_20260922
\par
La section~\ref{delta22:xi:holonomia} explicite la contribution de l'archive à la composition bilatérale, construit le plan de quart de tour du calendrier dodécaphasique et démontre la restitution orientée de l'action au moyen d'un commutateur relatif. L'amplification compatible conserve l'opérateur ; le détecteur distingue la dilatation de l'état du gain énergétique et l'archive ordonnée du décompte des opérations.
% END SINTESIS_ADITIVA_20260922

La réalisation de la section~\ref{xii:radial:section} compose l’inverseur
incidentiel avec l’inscription marquée et la norme locale d’action pour
obtenir le retour longitudinal \(L=L_*\alpha^{16}5759/23040\).
Son module d’arrivée conserve le caractère positif complet et détermine
un opérateur radial unique pour chaque caractère. Avec l’identification explicite
du rayon gravitationnel réduit, les deux réciprocités fixent le couplage
commun \(G_a=c^3L^2/\hbar_a\). La mémoire résiduelle, la variation conjointe
du champ, du transport et du poids, et la récupération des secteurs de bord
maintiennent leurs preuves contiguës. La section chronologique compatible transporte
la coordonnée électronique ; l’échelle surfacique \(L^2\) demeure commune aux
sections d’action. La restitution constitutive conserve en outre la collection
orientée dont l’évaluation à l’extrémité détermine la constante de Catalan.

% BEGIN R27_FRAMING_CONCLUSION_ES
L'antécédent fini de la récupération est déterminé plus précisément par les preuves de l'émetteur. Les deux signatures produisent \(42\) blocs possibles et le troisième chiffre satisfait une congruence de contrôle. Les raccordements régionaux conservent des témoins explicites dans les germes ; l'archive ordonnée distingue des registres qui partagent un histogramme et diffèrent par leur corrélation. Les jauges et les contrôles de clôture séparent la normalisation de l'observable des conditions qui fixent une évaluation. La frontière saturée admet huit distributions, la neutralité duale conserve une paire et l'orientation sélectionne son représentant. L'information que reçoit chaque lecteur est ainsi spécifiée avant l'application des identités bilatérales, des opérateurs d'inversion et des bornes de stabilité.
% END R27_FRAMING_CONCLUSION_ES

% BEGIN R28_FRAMING_CONCLUSION_ES
Les cinq clôtures régionales atteignent le minimum huit dans le problème fini de couverture des quatre arêtes prescrites. Le graphe conserve les étiquettes des émissions et le quotient conserve leurs incidences entre classes de phase. L'inversion orientée rend impair le lecteur \(\nu_{120}\) et maintient pair \(\nu_{270}\) ; ces identités explicitent la réponse de la mémoire au changement d'orientation.
% END R28_FRAMING_CONCLUSION_ES


La phase nonadique et le plan de quart de tour retrouvent exactement un
quotient de \(36\) positions. La proposition~\ref{ins29:xii:extension}
localise la fibre ternaire qui restitue les \(108\) positions et conserve
l’obstruction à scinder l’horloge. La mémoire de la dynamique enrichie
continue de distinguer ses retours complets.


Le résultat central est une caractérisation calculable de la conservation évolutive de l’unité : la structure détermine une évaluation, ses transformations distribuent l’information et les inverses précisent comment la retrouver. La constante de couplage, les lois des observables et les bornes de reconstruction expriment cette continuité au moyen d’équations et d’opérations, depuis le registre arithmétique jusqu’aux réalisations spécifiées.
